\uuid{mmKr}
\titre{Composée/DL usuels}
\auteur{hanani ; mzari}
\organisation{exo7}
\licence{CC BY-NC-SA 4.0}
\datecreate{2020-01-20}
\niveau{L1}
\module{Analyse}
\chapitre{Développement limité}
\sousChapitre{Calculs}
\difficulte{3}
\theme{Opérations sur les DL, Composée, DL usuels}
\exo7qcmid{lille-2.developpements-limites.018}

\contenu{
\texte{
Soit $a\in \Rr$ et $f$ une fonction. On dit que $f(x)$ est un $o\left((x-a)^n\right)$ au voisinage de $a$ et on écrit $f(x)=o((x-a)^n)$ si $\displaystyle \lim _{x\to a}\frac{f(x)}{(x-a)^n}=0$. On note par $DL_n(a)f(x)$ le développement limité de $f$ à l'ordre $n$ en $a$. On note aussi $DL_n(a^+)f(x)$ (resp. $DL_n(a^-)f(x)$) le développement limité de $f$ à l'ordre $n$ à droite en $a$ (resp. à gauche en $a$).
}

\begin{question}
\qclassification{125.02}
\qtags{motcle=Composée/DL usuels}
Soit $\displaystyle f(x)=\ln \left[1+\sin x\right]$ et $\displaystyle g(x)=\ln \left[1+\cos x\right]$. Parmi les affirmations suivantes, cocher celles qui sont vraies :
\begin{answers}
\good{$\displaystyle DL_2(0)f(x)=x-\frac{x^2}{2}+o(x^2)$.}
\bad{$\displaystyle DL_2(0)g(x)=1-\frac{x^2}{2}+o(x^2)$.}
\bad{$\displaystyle DL_2(0)\left[f(x)+g(x)\right]=1+x-x^2+o(x^2)$.}
\good{$\displaystyle DL_2(0)f(x)g(x)=\ln (2)x-\frac{\ln (2)}{2}x^2+o(x^2)$.}
\end{answers}
\vskip2mm
\begin{explanations}
Le $\displaystyle DL_2(0)\ln (1+t)=t-\frac{t^2}{2}+o(t^2)$ avec $\displaystyle t=\sin x=x+o(x^2)$ donne :
$$DL_2(0)f(x)=x-\frac{x^2}{2}+o(x^2.$$
Ensuite, on transforme $g(x)$ : $\displaystyle g(x)=\ln \left[2\left(1+\frac{\cos x-1}{2}\right)\right]=\ln (2)+\ln \left(1+\frac{\cos x-1}{2}\right)$. Ainsi, avec $\displaystyle t=\frac{\cos x-1}{2}=-\frac{x^2}{4}+o(x^2)$, on obtient : $\displaystyle g(x)=\ln (2)-\frac{x^2}{4}+o(x^2)$.
\end{explanations}
\end{question}
}
